\documentclass[10pt,a4paper]{article}

\usepackage[utf8]{inputenc}
\usepackage[T1]{fontenc}
\usepackage[french]{babel}
\usepackage[margin=1.5cm]{geometry}
\usepackage{amsmath,amssymb}
\usepackage{xcolor}
\usepackage{listings}
\usepackage{tcolorbox}
\tcbuselibrary{breakable}
\usepackage{enumitem}
\usepackage{array}

\lstset{
    language=C,
    basicstyle=\ttfamily\scriptsize,
    backgroundcolor=\color{gray!10},
    keywordstyle=\color{blue},
    commentstyle=\color{gray},
    stringstyle=\color{red},
    numbers=none,
    breaklines=true,
    frame=single,
    tabsize=4,
    showstringspaces=false,
    inputencoding=utf8,
    extendedchars=true,
    literate={é}{{\'e}}1 {è}{{\`e}}1 {à}{{\`a}}1 {ç}{{\c{c}}}1 
             {ê}{{\^e}}1 {î}{{\^i}}1 {ô}{{\^o}}1 {û}{{\^u}}1 {ù}{{\`u}}1
}

\newtcolorbox{infobox}[1]{colback=blue!5,colframe=blue!75!black,fonttitle=\bfseries,title=#1,breakable}
\newtcolorbox{examplebox}[1]{colback=green!5,colframe=green!50!black,fonttitle=\bfseries,title=#1,breakable}
\newtcolorbox{exercicebox}[1]{colback=orange!5,colframe=orange!75!black,fonttitle=\bfseries,title=#1,breakable}

\title{\textbf{Chapitre 2 : Organisation Mémoire}\\
\large Exercices d'Application (à faire pendant le cours)}
\author{Dr. El Hadji Bassirou TOURÉ\\
Département de Mathématiques et Informatique\\
Université Cheikh Anta Diop de Dakar}
\date{}

\begin{document}
\maketitle

\begin{infobox}{Objectif}
Ces exercices sont à réaliser \textbf{pendant le cours}, immédiatement après chaque notion. Ils permettent de vérifier la compréhension avant de passer aux TDs.
\end{infobox}

%==============================================================================
\section{Exercice A1 : Endianness (Little-Endian)}
%==============================================================================

\begin{exercicebox}{Énoncé}
La valeur \texttt{0xDEADBEEF} est stockée à l'adresse \texttt{0x100} sur une machine \textbf{little-endian}.

Remplir le tableau montrant le contenu de la mémoire :
\end{exercicebox}

\vspace{0.5em}

\begin{center}
\renewcommand{\arraystretch}{1.5}
\begin{tabular}{|c|c|c|c|c|}
\hline
\textbf{Adresse} & 0x100 & 0x101 & 0x102 & 0x103 \\
\hline
\textbf{Contenu (hex)} & \rule{1.5cm}{0.4pt} & \rule{1.5cm}{0.4pt} & \rule{1.5cm}{0.4pt} & \rule{1.5cm}{0.4pt} \\
\hline
\end{tabular}
\end{center}

\begin{examplebox}{Méthode}
\begin{enumerate}
    \item \textbf{Little-endian} = l'octet de poids faible (LSB) est stocké à l'adresse la plus basse
    \item Pour \texttt{0xDEADBEEF} : les octets sont DE, AD, BE, EF
    \item L'octet de poids faible est \texttt{EF}, il va à l'adresse 0x100
    \item Puis \texttt{BE} à 0x101, \texttt{AD} à 0x102, \texttt{DE} à 0x103
\end{enumerate}
\end{examplebox}

\vspace{1em}

\begin{exercicebox}{À vous !}
Sur la même machine little-endian, la valeur \texttt{0x12345678} est stockée à l'adresse \texttt{0x200}.

\begin{center}
\renewcommand{\arraystretch}{1.5}
\begin{tabular}{|c|c|c|c|c|}
\hline
\textbf{Adresse} & 0x200 & 0x201 & 0x202 & 0x203 \\
\hline
\textbf{Contenu} & \rule{1.5cm}{0.4pt} & \rule{1.5cm}{0.4pt} & \rule{1.5cm}{0.4pt} & \rule{1.5cm}{0.4pt} \\
\hline
\end{tabular}
\end{center}

Si on lit un \texttt{short} (2 octets) à l'adresse 0x200, quelle valeur obtient-on ? \rule{2cm}{0.4pt}
\end{exercicebox}

\newpage
%==============================================================================
\section{Exercice A2 : Pointeurs — Valeur et Adresse}
%==============================================================================

\begin{exercicebox}{Énoncé}
Soit le code suivant sur une machine 32 bits :
\begin{lstlisting}
int x = 42;       /* &x = 0x1000 */
int *p = &x;      /* &p = 0x2000 */
\end{lstlisting}

Compléter le tableau :
\end{exercicebox}

\vspace{0.5em}

\begin{center}
\renewcommand{\arraystretch}{1.5}
\begin{tabular}{|l|c|}
\hline
\textbf{Question} & \textbf{Réponse} \\
\hline
Valeur de \texttt{x} & \rule{3cm}{0.4pt} \\
\hline
Adresse de \texttt{x} (\texttt{\&x}) & \rule{3cm}{0.4pt} \\
\hline
Valeur de \texttt{p} & \rule{3cm}{0.4pt} \\
\hline
Adresse de \texttt{p} (\texttt{\&p}) & \rule{3cm}{0.4pt} \\
\hline
Valeur pointée \texttt{*p} & \rule{3cm}{0.4pt} \\
\hline
Taille de \texttt{p} (sizeof) & \rule{3cm}{0.4pt} \\
\hline
\end{tabular}
\end{center}

\begin{examplebox}{Méthode}
Un pointeur a \textbf{deux aspects} :
\begin{itemize}
    \item Sa \textbf{valeur} = l'adresse qu'il contient (vers où il pointe)
    \item Son \textbf{adresse} = où il est stocké en mémoire (\texttt{\&p})
\end{itemize}

L'opérateur \texttt{*} (déréférencement) permet d'accéder à la valeur pointée.
\end{examplebox}

\vspace{1em}

\begin{exercicebox}{À vous !}
Avec le même code, dessiner la mémoire aux adresses 0x1000 et 0x2000 (4 octets chacune, little-endian).

\vspace{0.5em}
\textbf{À 0x1000 (variable x = 42 = 0x2A) :}

\begin{center}
\begin{tabular}{|c|c|c|c|}
\hline
0x1000 & 0x1001 & 0x1002 & 0x1003 \\
\hline
\rule{1cm}{0.4pt} & \rule{1cm}{0.4pt} & \rule{1cm}{0.4pt} & \rule{1cm}{0.4pt} \\
\hline
\end{tabular}
\end{center}

\textbf{À 0x2000 (pointeur p = 0x1000) :}

\begin{center}
\begin{tabular}{|c|c|c|c|}
\hline
0x2000 & 0x2001 & 0x2002 & 0x2003 \\
\hline
\rule{1cm}{0.4pt} & \rule{1cm}{0.4pt} & \rule{1cm}{0.4pt} & \rule{1cm}{0.4pt} \\
\hline
\end{tabular}
\end{center}
\end{exercicebox}

\newpage
%==============================================================================
\section{Exercice A3 : Arithmétique des Pointeurs}
%==============================================================================

\begin{exercicebox}{Énoncé}
Soit le tableau suivant (\texttt{int} = 4 octets) :
\begin{lstlisting}
int arr[5] = {10, 20, 30, 40, 50};  /* &arr[0] = 0x1000 */
int *p = arr;
\end{lstlisting}

Calculer les adresses et valeurs :
\end{exercicebox}

\vspace{0.5em}

\begin{center}
\renewcommand{\arraystretch}{1.5}
\begin{tabular}{|l|c|c|}
\hline
\textbf{Expression} & \textbf{Adresse/Valeur} & \textbf{Calcul} \\
\hline
\texttt{\&arr[0]} & 0x1000 & (donné) \\
\hline
\texttt{\&arr[1]} & \rule{2cm}{0.4pt} & 0x1000 + \rule{1cm}{0.4pt} $\times$ 4 \\
\hline
\texttt{\&arr[3]} & \rule{2cm}{0.4pt} & \rule{3cm}{0.4pt} \\
\hline
\texttt{p + 2} & \rule{2cm}{0.4pt} & \rule{3cm}{0.4pt} \\
\hline
\texttt{*(p + 2)} & \rule{2cm}{0.4pt} & (valeur à cette adresse) \\
\hline
\texttt{p[4]} & \rule{2cm}{0.4pt} & équivalent à \rule{2cm}{0.4pt} \\
\hline
\end{tabular}
\end{center}

\begin{examplebox}{Formule clé}
\[\boxed{\texttt{p + n} = \texttt{p} + n \times \texttt{sizeof(*p)}}\]

Pour un \texttt{int*}, ajouter 1 au pointeur ajoute 4 octets à l'adresse !

\textbf{Équivalences :}
\begin{itemize}
    \item \texttt{arr[i]} $\equiv$ \texttt{*(arr + i)}
    \item \texttt{\&arr[i]} $\equiv$ \texttt{arr + i}
\end{itemize}
\end{examplebox}

\vspace{1em}

\begin{exercicebox}{À vous !}
Soit \texttt{short tab[4]} avec \texttt{\&tab[0] = 0x2000}. (\texttt{short} = 2 octets)

\begin{center}
\renewcommand{\arraystretch}{1.3}
\begin{tabular}{|l|c|}
\hline
\texttt{\&tab[2]} = & \rule{3cm}{0.4pt} \\
\hline
\texttt{(char*)tab + 5} = & \rule{3cm}{0.4pt} \\
\hline
\texttt{(int*)tab + 1} = & \rule{3cm}{0.4pt} \\
\hline
\end{tabular}
\end{center}
\end{exercicebox}

\newpage
%==============================================================================
\section{Exercice A4 : Chaînes — strlen vs sizeof}
%==============================================================================

\begin{exercicebox}{Énoncé}
Comparer les deux déclarations suivantes :
\begin{lstlisting}
char arr[] = "UCAD";
char *ptr = "UCAD";
\end{lstlisting}
\end{exercicebox}

\vspace{0.5em}

\begin{center}
\renewcommand{\arraystretch}{1.5}
\begin{tabular}{|l|c|c|}
\hline
\textbf{Expression} & \textbf{Pour arr} & \textbf{Pour ptr (64 bits)} \\
\hline
\texttt{sizeof(...)} & \rule{2cm}{0.4pt} & \rule{2cm}{0.4pt} \\
\hline
\texttt{strlen(...)} & \rule{2cm}{0.4pt} & \rule{2cm}{0.4pt} \\
\hline
Modification \texttt{[0]='X'} possible ? & \rule{2cm}{0.4pt} & \rule{2cm}{0.4pt} \\
\hline
Réaffectation \texttt{= "autre"} possible ? & \rule{2cm}{0.4pt} & \rule{2cm}{0.4pt} \\
\hline
\end{tabular}
\end{center}

\begin{examplebox}{Différences clés}
\begin{itemize}
    \item \texttt{char arr[]} : \textbf{tableau} alloué sur la pile, contenu \textbf{modifiable}
    \item \texttt{char *ptr} : \textbf{pointeur} vers une chaîne littérale en mémoire \textbf{read-only}
    \item \texttt{sizeof(arr)} = taille du tableau (inclut \texttt{'$\backslash$0'})
    \item \texttt{sizeof(ptr)} = taille du pointeur (4 ou 8 octets selon l'architecture)
    \item \texttt{strlen} = nombre de caractères \textbf{avant} le \texttt{'$\backslash$0'}
\end{itemize}
\end{examplebox}

\vspace{1em}

\begin{exercicebox}{À vous !}
Représenter la chaîne "UCAD" en mémoire (5 octets, avec le terminateur) :

\begin{center}
\renewcommand{\arraystretch}{1.5}
\begin{tabular}{|c|c|c|c|c|c|}
\hline
\textbf{Index} & 0 & 1 & 2 & 3 & 4 \\
\hline
\textbf{Caractère} & \rule{0.8cm}{0.4pt} & \rule{0.8cm}{0.4pt} & \rule{0.8cm}{0.4pt} & \rule{0.8cm}{0.4pt} & \rule{0.8cm}{0.4pt} \\
\hline
\textbf{Hex} & \rule{0.8cm}{0.4pt} & \rule{0.8cm}{0.4pt} & \rule{0.8cm}{0.4pt} & \rule{0.8cm}{0.4pt} & \rule{0.8cm}{0.4pt} \\
\hline
\end{tabular}
\end{center}

\textit{Table ASCII : U=0x55, C=0x43, A=0x41, D=0x44, NUL=0x00}
\end{exercicebox}

%==============================================================================
\vspace{2em}
\begin{infobox}{Récapitulatif des formules}
\begin{itemize}
    \item \textbf{Arithmétique pointeurs :} $\texttt{p + n} = \texttt{p} + n \times \texttt{sizeof(*p)}$
    \item \textbf{Équivalences :} \texttt{p[i]} $\equiv$ \texttt{*(p + i)} et \texttt{\&p[i]} $\equiv$ \texttt{p + i}
    \item \textbf{Chaînes :} \texttt{strlen} compte sans le \texttt{'$\backslash$0'}, \texttt{sizeof} compte avec
    \item \textbf{Little-endian :} LSB à l'adresse basse
\end{itemize}
\end{infobox}

\end{document}
